% Shared style; keep topic content and figures in this deck.
\newcommand{\CourseTemplatePath}{../../../../../templates/slides/}
\input{\CourseTemplatePath course_preamble.tex}
\title[Markov Decision Processes]{Lecture: Introduction to Markov Decision Processes}
\subtitle{States, transitions, actions, and policies}
\begin{document}
\makecoursetitle


\begin{frame}{The Markov Process: Motivative Examples}
\begin{itemize}
  \item \courseexample{Example: Days of the Week}
    \[\mathcal S=\{\text{Mon},\ldots,\text{Sun}\}.\]
    Tomorrow follows deterministically from today. Starting on Monday, we return to Monday after 7, 14, 21, \ldots\ days.
    \par\textit{Idea: a Markov process can be deterministic and cyclical.}
  \item \courseexample{Example: An AR(1) Process}
    \[s_{t+1}=\rho s_t+\varepsilon_{t+1},\qquad \varepsilon_{t+1}\overset{\mathrm{i.i.d.}}\sim N(0,\sigma^2).\]
    Assume \(\modelvar{\sigma>0}\); shocks are independent of the initial state. The current value determines the distribution of the next value.
    \par\textit{Idea: states can be continuous.} The process is Markov for any fixed \(\modelvar{\rho}\); \(\modelvar{|\rho|<1}\) gives a stable AR(1) with a stationary distribution.
\end{itemize}
\end{frame}

\begin{frame}{The Markov Process: Motivative Examples}
\courseexample{Example: Employment and Unemployment}
\begin{itemize}
  \item Each month, an unemployed worker finds a job with probability \(\modelvar{\alpha}\); an employed worker loses a job with probability \(\modelvar{\delta}\), where \(\modelvar{0<\alpha,\delta<1}\).
  \item With states ordered as unemployed (U), employed (E),
    \[P=\begin{pmatrix}1-\alpha&\alpha\\\delta&1-\delta\end{pmatrix}.\]
    Rows give today's state; columns give next month's state. Each row sums to one.
  \item \textit{Idea: transitions are random, even after we observe the state.}
    These probabilities are taken as given; the worker chooses no action in this model.
  \item How long does unemployment last? What fraction of time is spent unemployed in the long run?
\end{itemize}
{\footnotesize Source: \href{https://intro.quantecon.org/markov_chains_I.html}{QuantEcon, Markov Chains: Basic Concepts}.}
\end{frame}

\begin{frame}{The Markov Process: Motivative Examples}
\small
\setlength{\abovedisplayskip}{3pt}\setlength{\belowdisplayskip}{3pt}
\courseexample{Example: U.S. Business Cycle Regimes}
\par Monthly transition probabilities estimated by Hamilton (2005), as presented in \href{https://intro.quantecon.org/markov_chains_I.html}{QuantEcon}:
\[P=\begin{pmatrix}0.971&0.029&0\\0.145&0.778&0.077\\0&0.508&0.492\end{pmatrix}.\]
States: normal growth (ng), mild recession (mr), severe recession (sr).
\par\textit{Idea: persistence differs across states.} The probability of staying one more month is 97.1\% in ng, 77.8\% in mr, and 49.2\% in sr.
\par\vspace{2pt}{\centering
\begin{tikzpicture}[scale=0.9,transform shape,font=\sffamily\small,
          regime/.style={ellipse,draw=SessionCharcoal,fill=white,line width=0.7pt,
            minimum width=1.2cm,minimum height=0.7cm},
          transition/.style={draw=SessionCharcoal,line width=0.8pt,
            -{Stealth[length=2mm,width=1.5mm]}}]
          \node[regime] (ng) at (0,0) {ng};
          \node[regime] (mr) at (2.8,0) {mr};
          \node[regime] (sr) at (5.6,0) {sr};
          \path[transition]
            (ng) edge[loop above,min distance=8mm] node[above] {0.971} (ng)
            (mr) edge[loop above,min distance=8mm] node[above] {0.778} (mr)
            (sr) edge[loop above,min distance=8mm] node[above] {0.492} (sr)
            (ng) edge node[above] {0.029} (mr)
            (mr) edge[bend left=28] node[below] {0.145} (ng)
            (mr) edge node[above] {0.077} (sr)
            (sr) edge[bend left=28] node[below] {0.508} (mr);
        \end{tikzpicture}
\par}
\end{frame}

\begin{frame}{The Markov Process: Formal Definition}
\small
Let \(\modelvar{(s_t)_{t\geq0}}\) be a sequence of random variables taking values in \(\modelvar{\mathcal S}\).
\begin{itemize}
  \item We say the sequence satisfies the \term{Markov property} if, for every \(\modelvar{t}\) and every measurable subset \(\modelvar{\bar{\mathcal S}\subseteq\mathcal S}\),
  \[\Pr(s_{t+1}\in\bar{\mathcal S}\mid s_0,\ldots,s_t)
    =\Pr(s_{t+1}\in\bar{\mathcal S}\mid s_t)
    =P_t(\bar{\mathcal S}\mid s_t),\]
  where the \term{transition kernel} \(\modelvar{P_t(\bar{\mathcal S}\mid s)}\) gives the probability of entering \(\modelvar{\bar{\mathcal S}}\) from state \(\modelvar{s}\) at time \(\modelvar{t}\).
  \par Such a sequence is a \term{Markov process}: the current state contains all information from the history needed to predict the next state.
  \item We say a Markov process is \term{time-homogeneous} if \(\modelvar{P_t=P}\) for every \(\modelvar{t}\): its transition rule does not depend on calendar time.
  \item To specify its probability law, provide an \term{initial distribution} \(\modelvar{s_0\sim\mu_0}\) and the transition kernels \(\modelvar{(P_t)_{t\geq0}}\).
\end{itemize}
\end{frame}

\begin{frame}{The Markov Process: Transition Properties}
\small
Consider a time-homogeneous Markov process on a finite state space.
\begin{itemize}
  \item Its transition kernel is represented by a \term{transition matrix} \(\modelvar{P}\):
  \[P_{ij}=\Pr(s_{t+1}=j\mid s_t=i),\qquad P_{ij}\geq0,\qquad\sum_jP_{ij}=1.\]
  A matrix with nonnegative entries and unit row sums is called \term{row-stochastic}.
  \item If \(\modelvar{\mu_t}\) is a row vector of state probabilities,
  \[\mu_{t+1}=\mu_tP,\qquad \mu_{t+n}=\mu_tP^n.\]
  \item The probability of moving from \(\modelvar{i}\) to \(\modelvar{j}\) in \(\modelvar{n}\) steps is \(\modelvar{(P^n)_{ij}}\).
  Summing over intermediate states gives \(\modelvar{P^{n+m}=P^nP^m}\).
  \item \term{Conditional expectation:} for a column vector \(\modelvar{g}\) of state-dependent outcomes,
  \[\mathbb E[g(s_{t+1})\mid s_t=i]=\sum_jP_{ij}g(j)=(Pg)_i.\]
\end{itemize}
\end{frame}

\begin{frame}{The Markov Process: Reachability and Periodicity}
\small
Consider a finite, time-homogeneous chain.
\begin{itemize}
  \item We say the chain is \term{irreducible} if every state can reach every other state:
  \[\text{for every }i,j,\quad \exists n\geq0\text{ such that }(P^n)_{ij}>0.\]
  Positive one-step probabilities between every pair are unnecessary; a path through other states is enough.
  \item For a state \(\modelvar{i}\) with a possible return, define its \term{period} by
  \[d(i)=\gcd\{n\geq1:(P^n)_{ii}>0\}.\]
  An irreducible chain has a common period. We call it \term{aperiodic} if this period is 1, and \term{periodic} otherwise.
  \item The days-of-the-week chain is irreducible with period 7. The employment chain is irreducible and aperiodic: both states communicate and have self-loops.
\end{itemize}
\end{frame}

\begin{frame}{The Markov Process: Stationarity and Long-run Behavior}
\small
\setbeamertemplate{itemize/enumerate body begin}{\setlength{\itemsep}{1pt}}
\setlength{\abovedisplayskip}{4pt}\setlength{\belowdisplayskip}{4pt}
For a finite, time-homogeneous chain:
\begin{itemize}
  \item A probability row vector \(\modelvar{\pi}\) is a \term{stationary distribution} of the chain if
  \[\pi P=\pi,\qquad \pi_i\geq0,\qquad\sum_i\pi_i=1.\]
  Every finite chain has at least one. If \(\modelvar{\mu_0=\pi}\), every \(\modelvar{\mu_t=\pi}\).
  \item If the chain is \term{irreducible}, its stationary distribution is unique and
  \[\frac1T\sum_{t=0}^{T-1}\mathbf1\{s_t=i\}\longrightarrow\pi_i
    \quad\text{almost surely.}\]
  \item If the chain is also \term{aperiodic}, then for every initial distribution \(\modelvar{\mu_0}\),
  \[\mu_0P^n\longrightarrow\pi.\]
  The weekly cycle has uniform visit frequencies but its distribution keeps cycling when started on Monday.
\end{itemize}
\end{frame}

\begin{frame}{The Markov Decision Process: Motivative Examples}
\small
\courseexample{Example: Maze Escape}
\begin{columns}[c]
\begin{column}{0.58\textwidth}
\begin{itemize}
  \item \textbf{State:} current free cell. \textbf{Action:} up, down, left, or right.
  \item \textbf{Transition:} move one cell; hitting a wall or boundary leaves the position unchanged. The goal G is absorbing.
  \item \textbf{Reward:} \(\modelvar{-1}\) per move while outside G; zero thereafter. Choose actions to reach G quickly, with discount factor \(\modelvar{0<\beta<1}\).
\end{itemize}
\end{column}
\begin{column}{0.39\textwidth}
\begin{figure}
\centering
\begin{tikzpicture}[x=0.5cm,y=0.5cm,font=\sffamily\large]
      \definecolor{MazeWall}{HTML}{928580}
      \fill[white] (0,0) rectangle (9,6);
      \foreach \x/\y in {2/2,2/3,2/4,7/3,7/4,7/5,5/1}
        \fill[MazeWall] (\x,\y) rectangle +(1,1);
      \draw[SessionCharcoal,line width=0.3pt,step=1] (0,0) grid (9,6);
      \draw[SessionCharcoal,line width=0.9pt] (0,0) rectangle (9,6);
      \node[text=SessionCharcoal] at (0.5,3.5) {S};
      \node[text=SessionCharcoal] at (8.5,5.5) {G};
    \end{tikzpicture}
\caption{Maze navigation.}
\end{figure}
\end{column}
\end{columns}
\medskip
\textit{Idea: actions control the next state; rewards express the goal.}
\end{frame}

\begin{frame}{The Markov Decision Process: Motivative Examples}
\small
\courseexample{Example: McCall Job Search}
\begin{itemize}
  \item \textbf{State:} \(\modelvar{\mathrm U(w)}\), unemployed with current wage offer \(\modelvar{w}\), or \(\modelvar{\mathrm E(w)}\), employed at wage \(\modelvar{w}\).
  \item \textbf{Action and reward:} at \(\modelvar{\mathrm U(w)}\), accept and earn \(\modelvar{w}\) today, or reject and receive benefit \(\modelvar{b}\) today.
  \item \textbf{Transition:} rejection brings a new independent offer \(\modelvar{w'\sim F}\). Acceptance leads to \(\modelvar{\mathrm E(w)}\), with wage \(\modelvar{w}\) every period and no job separation. At \(\modelvar{\mathrm E(w)}\), the only action is to stay employed.
  \item \textbf{Objective:} maximize expected discounted income, \(\modelvar{0<\beta<1}\). Assume wages have bounded support and \(\modelvar{b}\) is finite.
  \item \textit{Idea: a discrete choice changes future opportunities.} Waiting preserves the chance of a better offer; accepting locks in income.
\end{itemize}
{\footnotesize Source: \href{https://python.quantecon.org/mccall_model.html}{QuantEcon, The McCall Search Model}.}
\end{frame}

\begin{frame}{The Markov Decision Process: Motivative Examples}
\courseexample{Example: Cake Eating}
\begin{itemize}
  \item \textbf{State:} remaining stock \(\modelvar{k_t\in[0,\bar k]}\), where \(\modelvar{\bar k<\infty}\).
  \item \textbf{Action:} consume \(\modelvar{c_t\in[0,k_t]}\).
  \item \textbf{Transition:} \(\modelvar{k_{t+1}=k_t-c_t}\); the unconsumed stock is saved. There is no production or uncertainty.
  \item \textbf{Reward:} \(\modelvar{u(c_t)=\sqrt{c_t}}\), including \(\modelvar{u(0)=0}\).
  \item \textbf{Objective:} choose consumption over time to maximize
  \[\sum_{t=0}^{\infty}\beta^t\sqrt{c_t},\qquad 0<\beta<1.\]
  \item \textit{Idea: both states and actions can be continuous, and feasible actions depend on the state.} Consuming more today leaves less for tomorrow.
\end{itemize}
\end{frame}

\begin{frame}{The Markov Decision Process: Formal Definition}
\small
We focus on a fully observed, time-homogeneous, infinite-horizon discounted MDP:
\[(\mathcal S,\{\mathcal A(s)\}_{s\in\mathcal S},P,r,\beta),\qquad 0<\beta<1.\]
\begin{itemize}
  \item \textbf{States and actions:} observe \(\modelvar{s_t\in\mathcal S}\), then choose \(\modelvar{a_t\in\mathcal A(s_t)}\), a nonempty feasible set.
  \item The state dynamics satisfy the \term{controlled Markov property} if, for every \(\modelvar{t}\) and measurable \(\modelvar{\bar{\mathcal S}\subseteq\mathcal S}\),
  \[\Pr(s_{t+1}\in\bar{\mathcal S}\mid s_0,a_0,\ldots,s_t,a_t)=P(\bar{\mathcal S}\mid s_t,a_t).\]
  where \(\modelvar{P}\) is the time-independent \term{controlled transition kernel}.
  \item The \term{reward function} \(\modelvar{r}\) specifies expected payoffs: receive \(\modelvar{R_{t+1}}\) after action \(\modelvar{a_t}\), with conditional mean \(\modelvar{r(s_t,a_t)}\) given the history and current action. Rewards and next states may be jointly random.
  \item The \term{discount factor} \(\modelvar{\beta}\) weights future rewards. Time homogeneity means that both \(\modelvar{P}\) and \(\modelvar{r}\) are independent of \(\modelvar{t}\).
\end{itemize}
\end{frame}

\begin{frame}{The Markov Decision Process: Timing}
\begin{figure}
    \centering
    % Vector redraw preserving Sutton and Barto's Fig. 3.1 topology.
    % Text inherits the same Source Sans 3 face as the slide body.
    \begin{tikzpicture}[x=1cm,y=1cm,font=\sffamily\normalsize,
      block/.style={draw=SessionCharcoal,line width=0.9pt,rounded corners=2mm,
        fill=white,minimum height=1cm,inner sep=3mm,font=\sffamily\Large},
      flow/.style={draw=SessionCharcoal,line width=1.1pt,-{Stealth[length=2.8mm,width=2mm]}},
      wire/.style={draw=SessionCharcoal,line width=1.1pt}]
      \node[block,minimum width=2.2cm] (agent) at (0,1.4) {Agent};
      \node[block,minimum width=4.3cm,minimum height=1.25cm] (environment) at (0,-1) {Environment};
      % Action: agent -> environment, on the right.
      \draw[flow] (agent.east) -- (4.2,1.4) -- (4.2,-1) -- (environment.east);
      \node[align=left,anchor=west] at (4.3,0.25) {action\\$A_t$};
      % State: environment -> next-state label -> outer return path -> agent.
      \draw[flow] ($(environment.west)+(0,-0.35)$) -- (-2.9,-1.35);
      \draw[wire] (-2.9,-1.35) -- (-4.4,-1.35) -- (-4.4,1.62);
      \draw[flow] (-4.4,1.62) -- ($(agent.west)+(0,0.22)$);
      \node[align=left,anchor=east] at (-4.55,0.2) {state\\$S_t$};
      \node[anchor=south,font=\sffamily\small] at (-2.55,-1.31) {$S_{t+1}$};
      % Reward: the parallel inner return path.
      \draw[flow] ($(environment.west)+(0,0.35)$) -- (-2.9,-0.65);
      \draw[wire] (-2.9,-0.65) -- (-3.8,-0.65) -- (-3.8,1.18);
      \draw[flow] (-3.8,1.18) -- ($(agent.west)+(0,-0.22)$);
      \node[align=left,anchor=west] at (-3.7,0.2) {reward\\$R_t$};
      \node[anchor=south,font=\sffamily\small] at (-2.55,-0.61) {$R_{t+1}$};
      \draw[SessionMuted,densely dashed,line width=0.45pt] (-2.9,-1.8) -- (-2.9,-0.25);
    \end{tikzpicture}
    \caption{\footnotesize Agent–environment loop (adapted from \href{http://incompleteideas.net/book/the-book-2nd.html}{Sutton and Barto (2018)}, Fig. 3.1).}
    \label{fig:sutton-mdp-loop}
  \end{figure}
\centering\small
Observe \(\modelvar{s_t}\), choose \(\modelvar{a_t}\), then receive \(\modelvar{R_{t+1}}\) and observe \(\modelvar{s_{t+1}}\).\\
\(\modelvar{r(s_t,a_t)}\) is the expected reward from that action.
\end{frame}

\begin{frame}{The Markov Decision Process: Policies}
\small
\begin{itemize}
  \item A policy specifies how actions are chosen. It is a \term{Markov policy} if its action distribution depends only on the current state and time.
  \item A Markov policy is \term{stationary} if it does not depend on time, and \term{deterministic} if it selects one action with probability one. Such a policy is a feasible rule
  \[a_t=\sigma(s_t),\qquad \sigma(s)\in\mathcal A(s).\]
  A stationary randomized policy instead specifies a distribution \(\modelvar{\sigma(\cdot\mid s)}\) over feasible actions.
\end{itemize}
\end{frame}

\begin{frame}{The Markov Decision Process: A Policy Induces a Chain}
\small
For a finite state space, fix a stationary deterministic policy \(\modelvar{\sigma}\).
\begin{itemize}
  \item The induced transition matrix and expected reward vector are
  \[P^\sigma_{ij}=P(j\mid i,\sigma(i)),\qquad r^\sigma_i=r(i,\sigma(i)).\]
  \item \(\modelvar{P^\sigma}\) is row-stochastic: fixing the decision rule gives a time-homogeneous Markov chain. Thus \(\modelvar{\mu_{t+1}=\mu_tP^\sigma}\).
  \item Different policies can change reachability and long-run behavior. An absorbing goal or accepted job can make \(\modelvar{P^\sigma}\) reducible; the earlier convergence conditions must be checked for each policy.
\end{itemize}
\end{frame}


% BEGIN MANAGED READINGS: mdp
\begin{frame}{References and Further Reading}
\coursereading{1}
  {Richard S. Sutton and Andrew G. Barto (2018), Reinforcement Learning: An Introduction. 2nd ed. MIT Press. [Book]}
  {https://www.andrew.cmu.edu/course/10-703/textbook/BartoSutton.pdf}
  {\begin{itemize}
\item \href{https://www.andrew.cmu.edu/course/10-703/textbook/BartoSutton.pdf\#page=69}{Chapter 3: Finite Markov Decision Processes}
\end{itemize}}
\coursereading{2}
  {Dimitri P. Bertsekas (2019), Reinforcement Learning and Optimal Control. Athena Scientific. [Book]}
  {https://faculty.engineering.asu.edu/bertsekas/books/reinforcement-learning-and-optimal-control}
  {\begin{itemize}
\item \href{https://faculty.engineering.asu.edu/sites/g/files/litvpz5506/files/2026-06/RL_1-SHORT-INTERNET-POSTED.pdf}{Chapter 1: Exact Dynamic Programming (selected sections)}
\end{itemize}}
\coursereading{3}
  {Dimitri P. Bertsekas and Yuchao Li (2025), Topics in Reinforcement Learning. CSE 691, Arizona State University. [Online course]}
  {https://web.mit.edu/dimitrib/www/RLbook.html}
  {\begin{itemize}
\item \href{https://web.mit.edu/dimitrib/www/RLTopics_2025_Lect1.pdf}{Lecture 1: Introduction to Sequential Decision Making (Reinforcement Learning, Dynamic Programming, Decision and Control)}
\item \href{https://web.mit.edu/dimitrib/www/RLTopics_2025_Lect2.pdf}{Lecture 2: Stochastic Finite and Infinite Horizon DP, Approximation in Value Space}
\end{itemize}}
\end{frame}
\begin{frame}{References and Further Reading}
\coursereading{4}
  {Thomas J. Sargent and John Stachurski, QuantEcon. [Online lectures]}
  {https://quantecon.org/lectures/}
  {\begin{itemize}
\item \href{https://intro.quantecon.org/markov_chains_I.html}{Markov Chains: Basic Concepts}
\item \href{https://intro.quantecon.org/markov_chains_II.html}{Markov Chains: Irreducibility and Ergodicity}
\item \href{https://python.quantecon.org/mccall_model.html}{Job Search I: The McCall Search Model}
\item \href{https://tools-techniques.quantecon.org/discrete_dp.html}{Discrete State Dynamic Programming}
\end{itemize}}
\end{frame}
% END MANAGED READINGS: mdp
\end{document}
